
\documentclass[../all/math-standard_questions_all.tex]{subfiles}

\begin{document}

\setlength{\abovedisplayskip}{2mm}
\setlength{\belowdisplayskip}{1mm}

\section{数学Ⅱ～複素数と方程式～}

\subsection{複素数［大阪薬科大 2006］}

虚数単位$i$に対して，
$i+i^2+i^3+ \dots +i^{50}$を計算せよ．

\vspace{2cm}
\vfill

\subsection{複素数［摂南大 2001］}

実数$x$，$y$が等式$(1+2i)x^2+(2+yi)x-3(1+i)=0$ ($i$は虚数単位)を満たすとき，
$x$，$y$の値を求めよ．

\vspace{2cm}
\vfill

\subsection{複素数［岡山理科大 2011］}

$k$を0と異なる実数の定数とし，$i$を虚数単位とする．
等式$x^2+(3+2i)x+k(2+i)^2=0$を満たす実数$x$が1つ存在するとし，それを$\alpha$とおくとき，
次の問いに答えよ．

\begin{enumerate}
    \item $k$と$\alpha$の値を求めよ．
    \item この等式を満たす複素数$x$をすべて求めよ．
\end{enumerate}

\vspace{2cm}
\vfill

\subsection{複素数［工学院大 2017］}

$x$の方程式$(a+i)x^2+2(1+i)x+1+ai=0$ ($i$は虚数単位)が
少なくとも1つの実数解をもつような実数$a$の値をすべて求めよ．
また，そのときの方程式の解をすべて求めよ．

\vspace{2cm}
\vfill

\subsection{複素数［東北大 2011］}

$a$を実数，$z$を0でない複素数とする．$z$と共役な複素数を$\overline{z}$で表す．

\begin{enumerate}

    \item $\displaystyle z+1-\frac{a}{\,z\,} = 0$を満たす$z$を求めよ．

    \item $\displaystyle \overline{z}+1-\frac{a}{\,\overline{z}\,} = 0$
        を満たす$z$が存在するような$a$の値の範囲を求めよ．

    \item $\displaystyle z(\overline{z})^2 + \overline{z} - \frac{a}{\,z\,} = 0$
        を満たす$z$が存在するような$a$の値の範囲を求めよ．

\end{enumerate}

\vspace{2cm}
\vfill

\subsection{解と係数の関係［大阪教育大 2023］}

2次方程式$x^2+x-1=0$の2つの解を$\alpha$，$\beta$とする．
次の式の値を求めよ．

\begin{enumerate}
    \item $(\alpha-1)(\beta-1)$
    \item $\alpha^4+\beta^4$
    \item $\alpha^{16}+\beta^{16}$
    \item $(\alpha+1)^8+(\beta+1)^8$
    \item $(\alpha^3+1)^8+(\beta^3+1)^8$
\end{enumerate}

\vspace{2cm}
\vfill

\subsection{解と係数の関係［関西大 1991］}

2次方程式$x^2+2px+p=0$の2つの解を$\alpha$，$\beta$とするとき，
$\displaystyle \frac{\alpha^2}{\beta}+ \frac{\beta^2}{\alpha}=-9$
となる定数$p$の値を求めよ．

\vspace{2cm}
\vfill

\subsection{解と係数の関係［常磐大 1991］}

2次方程式$x^2-(m+1)x+m=0$の1つの解が他の解の2倍に等しいという．
このときの$m$の値を求めよ．

\vspace{2cm}
\vfill

\subsection{解と係数の関係［法政大 2010］}

$a$，$b$を実数とする．2次方程式$x^2-ax+b=0$は2つの虚数解$\alpha$，$\beta$をもち，
$x^2+3ax+2b=0$の解は$\alpha^2$，$\beta^2$であるとする．
このとき，$a$および$b$を求めよ．

\vspace{2cm}
\vfill

\subsection{解と係数の関係［島根大 2011］}

$a$を実数とする．2次方程式$x^2+2ax+(a-1)=0$の解を$\alpha$，$\beta$とする．

\begin{enumerate}
    \item $\alpha$と$\beta$は異なる実数であることを示せ．
    \item $\alpha$と$\beta$のうち，少なくとも1つは負であることを示せ．
    \item $\alpha\leqq0$，$\beta\leqq0$であるとき，$\alpha^2+\beta^2$の最小値を求めよ．
\end{enumerate}

\vspace{2cm}
\vfill

\subsection{解と係数の関係［広島県立大 2004］}

実数$b$，$c$を係数とする2次方程式$x^2+bx+c=0$が異なる2つの虚数解をもつ．
1つの虚数解を$\alpha$とすると，他の解は$2\alpha-4+3i$と表すことができる．
このとき，$b$と$c$の値を求めよ．

\vspace{2cm}
\vfill

\subsection{剰余の定理と因数定理［浜松大 2004］}

整式$2x^4+7x^3+3x^2+2x-6$を$2x+1$で割ったときの商と余りを求めよ．

\vspace{2cm}
\vfill

\subsection{剰余の定理と因数定理［東京電機大 2016］}

$4x^{101}+3x^{100}-2x^{99}+1$を$x^3-x$で割った余りを求めよ．

\vspace{2cm}
\vfill

\subsection{剰余の定理と因数定理［立教大 2005］}

整式$P(x)$を$x+2$で割った余りが$3$，$x-3$で割った余りが$-1$のとき，
$P(x)$を$x^2- x - 6$で割った余りを求めよ．

\vspace{2cm}
\vfill

\subsection{剰余の定理と因数定理［岐阜聖徳学園大 2003］}

整式$f(x)$を$x^2+4x-5$で割った余りが$4x-3$，$2x^2-x-1$で割った余りが$5x-4$のとき，
$f(x)$を$x+5$，$2x^2+11x+5$で割った余りをそれぞれ求めよ．

\vspace{2cm}
\vfill

\subsection{剰余の定理と因数定理［千葉大 2004］}

整式$x^4+ax^3+ax^2+bx-6$が整式$x^2-2x+1$で割り切れるとき，$a$，$b$の値を求めよ．

\vspace{2cm}
\vfill

\subsection{剰余の定理と因数定理［名古屋学院大 2004］}

整式$f(x)$を$x^2+x+1$で割ると余りは$2x-1$であるという．

\begin{enumerate}

    \item $\left\{f(x)\right\}^2\left(=f(x)\times f(x)\right)$を
        $x^2+x+1$で割った余りを求めよ．
    \item $\left\{f(x)\right\}^2+af(x)+b$（ただし，$a$，$b$は実数の定数）
        が$x^2+x+1$で割り切れるときの$a$，$b$の値を求めよ．

\end{enumerate}

\vspace{2cm}
\vfill

\subsection{剰余の定理と因数定理［京都大 2019］}

$a$は実数とする．$x$に関する整式$x^5+2x^4+ax^3+3x^2+3x+2$を整式$x^3+x^2+x+1$で割ったときの
商を$Q(x)$，余りを$R(x)$とする．$R(x)$の1次の項の係数が1のとき，$a$の値を定め，
さらに$Q(x)$と$R(x)$を求めよ．

\vspace{2cm}
\vfill

\subsection{剰余の定理と因数定理［岩手大 2003］}

$x$の多項式$4x^3-2x^2-9x+7$を$x$の多項式$A$で割ると，その商が$B$で余りが$x+1$となる．
また，$A$と$B$の和は$2x^2+4x-5$である．このとき，$A$と$B$を求めよ．

\vspace{2cm}
\vfill

\subsection{剰余の定理と因数定理［神戸大 2012］}

$n$を自然数とし，多項式$P(x)$を$P(x)=(x+1)(x+2)^n$と定める．以下の問に答えよ．

\begin{enumerate}

    \item $P(x)$を$x-1$で割ったときの余りを求めよ．
    
    \item $(x+2)^n$を$x^2$で割ったときの余りを求めよ．
    
    \item $P(x)$を$x^2$で割ったときの余りを求めよ． 

    \item $P(x)$を$x^2(x-1)$で割ったときの余りを求めよ．

\end{enumerate}

\vspace{2cm}
\vfill

\subsection{高次方程式［学習院大 2005］}

$a$，$b$は実数で，方程式$x^3-2x^2+ax+b=0$は$x=2+i$を解にもつとする．
このとき，$a$，$b$の値と方程式のすべての解を求めよ．

\vspace{2cm}
\vfill

\subsection{高次方程式［高崎経済大 2006］}

実数解をもたない方程式$2x^4+x^3+x^2+x+2=0$を解きなさい．

\vspace{2cm}
\vfill

\subsection{高次方程式［自治医科大 2006］}

3次方程式$x^3-x^2+2x+3=0$の解を$\alpha$，$\beta$，$\gamma$とする．
$(\alpha+1)(\beta+1)(\gamma+1)-\alpha^2-\beta^2-\gamma^2$の値を求めよ．

\vspace{2cm}
\vfill

\subsection{高次方程式［岡山理科大 2014］}

正の数$x$，$y$，$z$が3条件
\begin{align*}
    \frac{\,1\,}{x}+\frac{\,1\,}{y}+\frac{\,1\,}{z}=\frac{\,7\,}{4}, \quad
    x^2+y^2+z^2=21, \quad xyz=8    
\end{align*}
を満たすとき，次の問いに答えよ．

\begin{enumerate}

    \item $xy+yz+zx$の値を求めよ．

    \item $x+y+z$の値を求めよ．
    
    \item $x\leqq y\leqq z$であるとき，$x$，$y$，$z$の値を求めよ．
    
\end{enumerate}

\vspace{2cm}
\vfill

\subsection{高次方程式［立命館大 2017］}

$a$，$b$，$c$を実数とする．3次方程式
\begin{align*}
    x^3+ax^2+bx+c=0 \myleader{1.8cm} \text{①}
\end{align*}
の解の1つが$-1+\sqrt{2}\,i$であるとき，次の問いに答えよ．

\begin{enumerate}

    \item $b$，$c$を$a$を用いて表せ．

    \item 方程式\,①\,の実数解を$a$を用いて表せ．
    
    \item 方程式\,①\,の実数解が，方程式$x^2+bx-b-1=0$の解でもあるとき，$a$の値を求めよ．
        ただし，$a>0$とする．
        \label{item3}

    \item $a$を\!\ref{item3}\!\!で求めた値とする．
        方程式\,①\,の3つの解を$\alpha$，$\beta$，$\gamma$とするとき，
        $\alpha^4+\beta^4+\gamma^4$の値を求めよ．

\end{enumerate}

\vspace{2cm}
\vfill

\subsection{高次方程式［お茶の水女子大 2008］}

次の問いに答えよ．

\begin{enumerate}

    \item $\displaystyle y=x+\frac{\,1\,}{x}$とおき，
        $\displaystyle x^3+\frac{1}{x^3}$，$\displaystyle x^4+\frac{1}{x^4}$を
        それぞれ$y$の多項式として表せ．

    \item $\alpha^6+\alpha^5-9\alpha^4-10\alpha^3-9\alpha^2+\alpha +1=0$を満たす
        すべての複素数$\alpha$を求めよ．

\end{enumerate}

\vspace{2cm}
\vfill

\subsection{高次方程式［大同工業大 2004］}

$8x^4+6x^2-9$を係数が有理数の範囲で因数分解せよ．
また，4次方程式$8x^4-6x^3+6x^2-9x-9=0$の実数解を求めよ．

\vspace{2cm}
\vfill

\subsection{高次方程式［横浜国立大 2014］}

正の定数$a$，$b$に対し$f(x)=ax^2-b$とおく．

\begin{enumerate}

    \item $f\left(f(x)\right)-x$は$f(x)-x$で割り切れることを示せ．

    \item 方程式 $f\left(f(x)\right)-x=0$が異なる4つの実数解をもつための
        $a$，$b$の条件を求めよ．

\end{enumerate}

\vspace{2cm}
\vfill

\subsection{高次方程式［お茶の水女子大 2002］}

\begin{enumerate}

    \item $(a+b+c)(a^2+b^2+c^2-ab-bc-ca)$を展開せよ．
        \label{item1}

    \item 関係式$\alpha\beta=12$，$\alpha^3+\beta^3=91$を満たす
        実数の組$(\alpha,\; \beta)$をすべて求めよ．
        \label{item2}

    \item 方程式$x^3-36x+91=0$の解をすべて求めよ．
        必要ならば，\!\!\ref{item1}\!\!\!\!，\!\!\!\ref{item2}\!\!\!を利用せよ．

\end{enumerate}

\vspace{2cm}
\vfill

\subsection{高次方程式［大阪工業大 2006］}

$a$を定数とする．$f(x)=ax^3-(a+2)x^2-(2a-5)x-2$について，次の問いに答えよ．

\begin{enumerate}

    \item $f(x)$を因数分解せよ．

    \item $f(x)=0$が重解をもつような$a$の値を求めよ．

\end{enumerate}

\vspace{2cm}
\vfill

\subsection{高次方程式［東北学院大 1991］}

2つの方程式
\begin{align*}
    x^3+px+2=0 \!\!\myleader{1cm} \text{①}\,, \quad
    3x^2+p=0 \!\!\myleader{1cm} \text{②}
\end{align*}
が共通解をもつとする．

\begin{enumerate}

    \item $p$の値を求めよ．

    \item このとき，①，②\,を解け．

\end{enumerate}

\vspace{2cm}
\vfill

\subsection{高次方程式［新潟大 2016］}

整式$P(x)=x^4+x^3+x-1$について，次の問いに答えよ．

\begin{enumerate}

    \item $i$を虚数単位とするとき，$P(i)$，$P(-i)$の値を求めよ．
    
    \item 方程式$P(x)=0$の実数解を求めよ．

    \item $Q(x)$を3次以下の整式とする．次の条件
    \begin{align*}
        & Q(1)=P(1), \quad Q(-1)=P(-1), \\
        & Q(2)=P(2), \quad Q(-2)=P(-2)
    \end{align*}
    をすべて満たす$Q(x)$を求めよ．

\end{enumerate}

\vspace{2cm}
\vfill


\end{document}